% One-time imported draft. Editable native source; no Markdown regeneration.
\section{Why the two error tests use these
thresholds}\label{007-bresenham-line-x-derivation-why-the-two-error-tests-use-these-thresholds}

Let A be the ordered start, dx \ensuremath{\geq} 0, D \ensuremath{\geq} 0, i = x\ensuremath{-}Ax, and j = sy\ensuremath{\times}(y\ensuremath{-}Ay).
The invariant is err = dx\ensuremath{-}D\ensuremath{-}D\ensuremath{\times}i+dx\ensuremath{\times}j. An x step increases i by one and
subtracts D; a y step increases j by one and adds dx.

On an x-major line, the next ideal minor coordinate is D\ensuremath{\times}(i+1)/dx. It
rounds one step forward when D\ensuremath{\times}(i+1)/dx\ensuremath{-}j \ensuremath{\geq} 1/2. Multiplying by the
positive denominator and substituting the invariant gives 2\ensuremath{\times}err \ensuremath{\leq} dx,
the y-step test.

On a y-major line, the next ideal minor coordinate is dx\ensuremath{\times}(j+1)/D. Its
forward-step condition dx\ensuremath{\times}(j+1)/D\ensuremath{-}i \ensuremath{\geq} 1/2 rearranges to 2\ensuremath{\times}err \ensuremath{\geq} \ensuremath{-}D, the
x-step test.

The non-major axis uses its rounding test; the major axis advances on
every update. The point case exits before a zero-denominator argument is
needed. Both tests use the same pre-update error. Equality chooses the
canonical forward minor step.

\subsection{Why the major axis
advances}\label{007-bresenham-line-x-derivation-why-the-major-axis-advances}

The nearest-cell rule also gives a bound that explains progress. On an
x-major line, write \ensuremath{\delta} = j \ensuremath{-} D\ensuremath{\times}i/dx. Nearest rounding gives
\textbar \ensuremath{\delta}\textbar{} \ensuremath{\leq} 1/2. Substituting into the error relation gives
err = dx \ensuremath{-} D + dx\ensuremath{\times}\ensuremath{\delta}, so err \ensuremath{\geq} dx/2 \ensuremath{-} D \ensuremath{\geq} \ensuremath{-}D/2 because dx \ensuremath{\geq} D. The x-step
condition therefore always succeeds on a nonfinal x-major sample.

On a y-major line, use \ensuremath{\delta} = i \ensuremath{-} dx\ensuremath{\times}j/D. The same rounding bound gives err
= dx \ensuremath{-} D \ensuremath{-} D\ensuremath{\times}\ensuremath{\delta} \ensuremath{\leq} dx \ensuremath{-} D/2 \ensuremath{\leq} dx/2 because D \ensuremath{\geq} dx. The y-step condition
always succeeds there. The other condition chooses whether the minor
coordinate steps too.

The ordered starting point satisfies the nearest-cell relation with i =
j = 0. The threshold calculation above gives the next rounded minor
coordinate, and the major-axis argument guarantees one unit of major
progress. Repeating that reasoning preserves the relation until the
endpoint. There are max(dx,D) updates and max(dx,D)+1 painted samples,
at most eight on this canvas. The endpoint test prevents an additional
update. The point case needs no update or division argument.

% Native mathematical source; requires amsmath. Integrate tags through the author template.
\begin{align*}
  i & = x-A_x,         & j & = s_y(y-A_y), \\
  e & = d_x-D-Di+d_xj.
\end{align*}
An $x$ step subtracts $D$ from $e$; a $y$ step adds $d_x$.
For an $x$-major line with $d_x>0$, the next minor coordinate steps forward when
\begin{align*}
  \frac{D(i+1)}{d_x}-j & \ge \frac12
                       & \Longleftrightarrow\quad 2e & \le d_x.
\end{align*}
For a $y$-major line with $D>0$, the corresponding condition is
\begin{align*}
  \frac{d_x(j+1)}{D}-i & \ge \frac12
                       & \Longleftrightarrow\quad 2e & \ge -D.
\end{align*}
Both decisions use the same pre-update value of $2e$. Equality chooses the forward
minor step after endpoint normalization. The coincident-point case exits before
these positive-denominator arguments are needed.


\paragraph{Major-axis progress.}
On an $x$-major line, let $\delta=j-Di/d_x$. Nearest rounding gives
$|\delta|\le\tfrac12$, hence
\[
  e=d_x-D+d_x\delta\ge\frac{d_x}{2}-D\ge-\frac D2
  \qquad(d_x\ge D).
\]
Thus the $x$ decision always succeeds before the endpoint.
On a $y$-major line, let $\delta=i-d_xj/D$. Then
\[
  e=d_x-D-D\delta\le d_x-\frac D2\le\frac{d_x}{2}
  \qquad(D\ge d_x),
\]
so the $y$ decision always succeeds. The minor-axis test gives the next rounded
sample. Starting at $i=j=0$, induction preserves that relation and reaches the
endpoint after $\max(d_x,D)$ updates. Painting before the endpoint test gives
$\max(d_x,D)+1$ samples, at most eight in this assignment. A point exits without
an update; each division in this argument is used only for a nonzero major extent.

